% !TeX program = xelatex
% Companion to the Lean sources in AIV5/mmrank21; see README.md in this directory.
\documentclass[11pt,a4paper]{article}
\usepackage[margin=25mm,headheight=15pt]{geometry}
\usepackage{fontspec}
\setmainfont{DejaVu Serif}[Scale=0.93]
\setsansfont{DejaVu Sans}[Scale=0.93]
\setmonofont{DejaVu Sans Mono}[Scale=0.8]
\usepackage[english]{babel}
\usepackage{amsmath,amssymb,amsthm,mathtools}
\usepackage{tikz-cd}
\usepackage{microtype}
\usepackage{enumitem}
\usepackage{needspace}
\usepackage{booktabs,longtable}
\usepackage[unicode,colorlinks=true,linkcolor=blue!45!black,urlcolor=blue!45!black]{hyperref}
\hypersetup{pdftitle={A lower bound of 21 for 3 x 3 matrix multiplication},
pdfauthor={Andrey Ivanov},
pdfsubject={Mathematical proof notes}}
\usepackage{fancyhdr}
\pagestyle{fancy}\fancyhf{}\fancyhead[L]{\small The lower bound $R_K\langle3,3,3\rangle\geqslant21$}\fancyhead[R]{\thepage}
\setlength{\parindent}{1.3em}\setlength{\parskip}{3pt}
\setlist{itemsep=3pt,topsep=5pt}
\emergencystretch=2em
% Keep inline formulas unbroken at binary operators and relations.
\binoppenalty=10000
\relpenalty=10000
\interdisplaylinepenalty=10000
\numberwithin{equation}{section}
\theoremstyle{plain}
\newtheorem{theorem}{Theorem}[section]
\newtheorem*{maintheorem}{Main theorem}
\newtheorem{lemma}[theorem]{Lemma}
\newtheorem{prop}[theorem]{Proposition}
\newtheorem{cor}[theorem]{Corollary}
\theoremstyle{definition}
\newtheorem{definition}[theorem]{Definition}
\newtheorem{remark}[theorem]{Remark}
\DeclareMathOperator{\tr}{tr}\DeclareMathOperator{\rk}{rk}
\DeclareMathOperator{\im}{im}\DeclareMathOperator{\supp}{supp}
\DeclareMathOperator{\Mat}{Mat}\DeclareMathOperator{\Span}{span}
\DeclareMathOperator{\Hom}{Hom}\DeclareMathOperator{\diag}{diag}
\newcommand{\kk}{K}\newcommand{\one}{\mathbf1}
\newcommand{\pair}[2]{\langle #1,#2\rangle}
\newcommand{\col}{\mathcal C}\newcommand{\test}{\mathcal T}
\newcommand{\rankone}[2]{#1#2}
\newcommand{\A}{\mathcal A}\newcommand{\B}{\mathcal B}\newcommand{\N}{\mathcal N}
\newcommand{\E}{\mathcal E}\newcommand{\Sig}{\Sigma}
\newcommand{\leqs}{\leqslant}\newcommand{\geqs}{\geqslant}
\begin{document}
\begin{center}
{\LARGE A lower bound for matrix multiplication}\\[6pt]
{\Large $R_K\langle3,3,3\rangle\geqslant21$}\\[8pt]
{\normalsize Mathematical proof notes}
\end{center}
\begin{quote}\small
These notes are an AI-generated, human-readable exposition of the Lean
proof in \href{https://github.com/AIV5/mmrank21}{AIV5/mmrank21}, not a
research article. They have not undergone independent mathematical or
peer review. No claim of novelty or priority is made for the results
or proofs. The Lean source is the authoritative formal statement and
proof; checking it does not certify the correctness of this prose account.
The accompanying source map links the statements to Lean modules.
\end{quote}
\begin{maintheorem}
Let $K$ be an arbitrary field. If the identity
\[
 BA=\sum_{i=1}^{r} a_i(A)b_i(B)W_i
\]
holds for all matrices $A,B\in\Mat_3(K)$, where $a_i,b_i$ are linear
forms and the matrices $W_i\in\Mat_3(K)$ are fixed, then $r\geqs21$.
\end{maintheorem}

\tableofcontents
\newpage
\section{Matrix notation and basic identities}

All vectors in $P,Q,S$ are regarded as row vectors; matrices act on the
right. For now we allow dimensions $p=\dim P$, $q=\dim Q$, $s=\dim S$,
where $p,s\geqs1$ and $q\geqs2$.
The coefficient matrices and inputs act in the following directions:
\[
\begin{tikzcd}[column sep=large]
 P\arrow[r,"U_i"] & Q\arrow[r,"V_i"] & S\arrow[ll,bend left=35,"W_i"]
\end{tikzcd}
\qquad
\begin{tikzcd}[column sep=large]
 S\arrow[r,"B"]&Q\arrow[r,"A"]&P\arrow[ll,bend left=35,"C"]
\end{tikzcd}
\]
For example, $U_i$ has size $p\times q$, $A$ has size $q\times p$,
and $BA$ has size $s\times p$. In a product $UV$, $U$ acts first and
$V$ second. A column $\eta$ defines the linear form $x\mapsto x\eta$
on row vectors. If $a$ is a nonzero column and $b$ a nonzero row,
the ordinary matrix product $ab$ has rank one.
We call the lines $Ka$ and $Kb$ its \emph{column direction} and
\emph{row direction}, respectively.
The \emph{column support} of a family of matrices is the span of all
their columns; its \emph{row support} is the span of all their rows,
equivalently the sum of their images under the right action.

\begin{definition}
A \emph{bilinear algorithm of length $r$} consists of $r$ triples of
matrices $(U_i,V_i,W_i)$, indexed by a set $I$, satisfying
\begin{equation}\label{eq:algorithm}
 BA=\sum_{i\in I}\tr(U_iA)\tr(V_iB)W_i\qquad\text{for all }A,B.
\end{equation}
All algorithms considered here are exact: the identity holds over the
field itself, with no limiting process. We omit this qualification below.
The minimum length for this format is denoted by
$R_K\langle p,q,s\rangle$.
\end{definition}
Every linear form on matrices is uniquely of the form
$A\mapsto\tr(UA)$: its value on the matrix unit $E_{ab}$ is $U_{ba}$.
Thus this definition agrees with the formulation of the main theorem.

\begin{definition}
A matrix $C:P\to S$, representing the output linear form
$BA\mapsto\tr(CBA)$, will be called a \emph{test}.
The \emph{maps defined by the coefficient forms} are denoted by
\[
 \alpha(A)_i=\tr(U_iA),\quad
 \beta(B)_i=\tr(V_iB),\quad
 \gamma(C)_i=\tr(W_iC).
\]
They take values in $K^I$. On this space we use the coordinatewise
product $xy=(x_i y_i)_i$, the identity $\one=(1)_i$, and the pairing
$\pair{x}{y}=\sum_i x_i y_i$. For subspaces $X,Y\subseteq K^I$,
$XY$ denotes the \emph{linear span} of all products $xy$.
The support of a vector $x$ is $\supp x=\{i:x_i\ne0\}$; the support
of a subspace is the union of the supports of its vectors.
\end{definition}
With respect to a nondegenerate pairing, the \emph{annihilator} of a
subspace consists of all elements of the other space that pair to zero
with it. For matrices we use $\tr(XY)$, and for coordinate vectors
we use $\pair{x}{y}$. In the latter case we also write $X^\perp$ and
speak of orthogonality; no positive definiteness is assumed.

\begin{lemma}\label{lem:identities}
The maps $\alpha,\beta,\gamma$ are injective. For all matrices of the
appropriate sizes, the following identities hold:
\begin{align}
 \tr(CBA)&=\sum_i\alpha(A)_i\beta(B)_i\gamma(C)_i,\label{eq:trace}\\
 CB&=\sum_i\beta(B)_i\gamma(C)_iU_i,&
 AC&=\sum_i\gamma(C)_i\alpha(A)_iV_i,\label{eq:cycles}\\
 \sum_i\gamma(C)_i U_i X V_i&=\tr(X)C
       &&(X\in\Mat_q(K)),\label{eq:sandwich}\\
 \sum_i\tr(XV_iTU_i)W_i&=\tr(X)T
       &&(T\in\Mat_{s,p}(K)).\label{eq:dual-sandwich}
\end{align}
\end{lemma}
\begin{proof}
The first identity follows by multiplying \eqref{eq:algorithm} by $C$
and taking the trace. Nondegeneracy of $\tr(XY)$, checked on matrix
units, gives both identities in \eqref{eq:cycles}.
If, for example, $\alpha(A)=0$, then $BA=0$ for every $B$, hence $A=0$.
Cyclic permutation gives the other two injectivity statements.

To prove \eqref{eq:sandwich}, first substitute matrix units for $A,B$
in \eqref{eq:trace}. Comparing coefficients gives
\[
 \sum_i(U_i)_{a b}(V_i)_{c d}\gamma(C)_i
       =\delta_{bc}C_{ad}.
\]
Multiplying by $X_{bc}$ and summing yields \eqref{eq:sandwich}.
Trace duality in $C$ gives \eqref{eq:dual-sandwich}.
All these identities are polynomial in the coefficients and remain
valid after a field extension.
\end{proof}

Zero summands may be deleted. Assuming an algorithm of length at most
twenty exists, pass to an algebraic closure of $K$ and delete its zero
summands. Unless stated otherwise, from this point until the final
section, $K$ is algebraically closed, all $U_i,V_i,W_i$ are nonzero,
and $r\leqs20$.

\section{Coordinatewise products and conductors}

\begin{definition}
For subspaces $M,N\subseteq K^I$, their \emph{conductor} is
\[
 (N:M)=\{d\in K^I:md\in N\text{ for every }m\in M\}.
\]
This is a subspace of $K^I$; the colon does not denote a quotient space.
The \emph{multiplicative stabilizer} of $X$ is the algebra
$(X:X)=\{a:aX\subseteq X\}$. Its elements need not be invertible.
\end{definition}

\begin{lemma}\label{lem:nz-growth}
Let $X,Y\subseteq K^I$ be nonzero subspaces such that $xy=0$ with
$x\in X$, $y\in Y$ implies $x=0$ or $y=0$. Then
\[
 \dim XY\geqs\dim X+\dim Y-1.
\]
\end{lemma}
\begin{proof}
Induct on $\dim X$. If $\dim X=1$, multiplication by any nonzero
$x\in X$ is injective on $Y$. Otherwise choose a coordinate $i$
on which some product is nonzero. The spaces
$X_0=\{x\in X:x_i=0\}$ and $(XY)_0=\{v\in XY:v_i=0\}$
have codimension one in $X$ and $XY$, respectively, and
$X_0Y\subseteq(XY)_0$. The induction hypothesis applied to $X_0,Y$
gives the required bound.
\end{proof}

\begin{definition}
For $G\subseteq I$, set
\[
 X[G]=\{x\in X:x_i=0\text{ if }i\notin G\},\qquad
 X|_G=\{(x_i)_{i\in G}:x\in X\}.
\]
These are the \emph{coordinate section} and the \emph{restriction} of
the space, respectively. Their dimensions need not be equal.
\end{definition}

\begin{lemma}\label{lem:two-conductors}
Let $M=\Span(\one,x,y)$ and $N,B\subseteq K^I$, with $MB\perp N$.
Then
\[
 \dim(N:M)+\dim(B:M)\geqs3\dim N+3\dim B-2r.
\]
\end{lemma}
\begin{proof}
For any $L\subseteq K^I$, consider
\[
 \Phi_L:L^3\longrightarrow K^I\oplus K^I,\qquad
 (l_0,l_1,l_2)\longmapsto(l_1-xl_0,l_2-yl_0).
\]
Its kernel consists of the triples $(l_0,xl_0,yl_0)$ with $l_0\in(L:M)$.
Thus
\[
 \dim\im\Phi_L=3\dim L-\dim(L:M).
\]
The form on $K^I\oplus K^I$ defined by
\[
 \omega((a,b),(c,d))=\pair{a}{d}-\pair{b}{c}
\]
is nondegenerate and alternating, meaning $\omega(v,v)=0$, including
in characteristic two. The images of $\Phi_N$ and $\Phi_B$ are
orthogonal: on expanding, the terms with factors $1,x,y$ vanish by
hypothesis, while the terms with $xy$ cancel. The dimensions of
orthogonal subspaces sum to at most $2r$. Substitute the dimensions
computed above.
\end{proof}

\subsection{Connectedness and coordinate components}

\begin{definition}
A finite nonempty family of nonzero vectors $(v_i)_{i\in G}$ is
\emph{connected} if every partition $G=G_1\sqcup G_2$ into two nonempty
parts satisfies
\[
 \Span\{v_i:i\in G_1\}\cap\Span\{v_i:i\in G_2\}\ne0.
\]
By splitting disconnected families, we obtain the unique finest
partition for which the span of the family is the direct sum of the
spans of the parts. These parts are called its connected components.
\end{definition}
Uniqueness also follows from a description by minimal linear
dependences: all indices in one such dependence must belong to the
same part; the transitive closure of these relations gives the components.

\begin{lemma}\label{lem:connected}
Suppose that the nonzero rank-one matrices $a_i b_i$ form a connected
family. Then
\[
 \dim\Span\{a_i\}+\dim\Span\{b_i\}
       \leqs\dim\Span\{a_i b_i\}+1.
\]
\end{lemma}
\begin{proof}
Start with one matrix. Suppose the inequality has been proved for the
part $J$ already chosen. If $J\ne G$, connectedness allows us to choose
a minimal nonempty part $H\subseteq G\setminus J$ whose span intersects
the span of $J$ nontrivially. The matrices indexed by $H$ are independent;
the intersection is one-dimensional and spanned by
$T=\sum_{i\in H}c_i a_i b_i\ne0$, where every $c_i\ne0$.
Otherwise one could remove an index from $H$.

Write $A_0,B_0$ for the previous spans of the factors and $A_1,B_1$
for the enlarged spans. The columns of $T$ lie in $A_0$, and its rows
lie in $B_0$. This implies
\[
 (\dim A_1-\dim A_0)+(\dim B_1-\dim B_0)\leqs |H|-1.
\]
Here is a direct linear algebra verification of this last step.
Factor $T=XY$, with columns $c_i a_i$ in $X$ and rows $b_i$ in $Y$.
Projecting the columns modulo $A_0$ gives $\bar X Y=0$; projecting
the rows modulo $B_0$ gives $X\bar Y=0$.
If the two dimension increments summed to $|H|$, we would have
$\im\bar Y=\ker\bar X$; then $X\bar Y=0$ would say that $X$
annihilates $\ker\bar X$. Since $\bar X Y=0$, this would imply
$XY=0$, contradicting $T\ne0$. The usual rank inequality for a product
first bounds the sum of the increments by $|H|$, leaving the bound
$|H|-1$. Adjoining $H$ increases the span of the matrices by exactly
$|H|-1$. The inequality is preserved, and finite induction completes
the proof.
\end{proof}

\begin{definition}
The \emph{coordinate components} of a subspace $X\subseteq K^I$ are
the connected components of its nonzero coordinate forms $x\mapsto x_i$.
Zero coordinate forms are omitted. Equivalently, these are the minimal
nonempty blocks $G$ for which $X=\bigoplus_G X[G]$.
We call these components \emph{active}: the space $X$ is not identically
zero on any coordinate of such a component.
\end{definition}

\begin{lemma}\label{lem:component-cost}
Let $\one\in M$, $D=(N:M)$, and $E=MD$. For each active coordinate
component $G$ of $E$, we have
\[
 \dim(M|_G)+\dim D[G]\leqs\dim E[G]+1.                 \tag{*}
\]
If $\dim E=\dim D+1$, there is such a component with $\dim(M|_G)=2$.
If $E=D$, then $D$ is spanned by common eigenvectors of all the
operators $w\mapsto mw$, $m\in M$.
\end{lemma}
\begin{proof}
Coordinate projections preserving $E$ also preserve $D$, since
\[
 m\pi_G(d)=\pi_G(md)\in E\subseteq N.
\]
Hence
$E[G]=(M|_G)D[G]$. Choose bases in $M|_G$ and $D[G]$.
The product coordinates are given by the bilinear forms
$(m,d)\mapsto m_i d_i$, whose matrices have rank one.
The linear relations among these matrices are exactly the relations
among the coordinate forms on $E[G]$, since the products span $E[G]$.
Thus the matrix family is connected, and its span has dimension
$\dim E[G]$. The first and second factors span spaces of dimensions
$\dim(M|_G)$ and $\dim D[G]$: they are the full sets of coordinate
forms on the corresponding spaces. Lemma~\ref{lem:connected} gives (*).

The sum of the dimension increments over the components is the total
increment. When the total increment is one, (*) gives
$\dim(M|_G)\leqs2$ for every component. If all these dimensions were one,
each $m$ would be constant on each block, because $\one\in M$.
Then $ME=E$, so $E\subseteq(N:M)=D$, a contradiction.

When $MD=D$, we have $M\subseteq(D:D)$: every multiplier belongs to
the multiplicative stabilizer of $D$.
The restrictions to $D$ of the diagonal operators $w\mapsto mw$ are
simultaneously diagonalizable. This can be seen directly: partition
$I$ according to equality of all the values $m_i$. By interpolation,
the projection onto each class is a polynomial in finitely many of
these diagonal operators, so it preserves $D$. On each resulting
coordinate section all the operators are scalar. Every nonzero vector
in that section is therefore a common eigenvector.
\end{proof}

\section{Normalizing the algorithm and dimension bounds}

\subsection{Choosing a hyperplane}
For a matrix subspace $H\subseteq\Mat_{p,q}(K)$, set
\[
 \col(H)=\{a\in K^{p\times1}:ab\in H
                 \text{ for all }b\in K^{1\times q}\}.
\]
This is a space of columns, not a kernel of the matrices in $H$.

\begin{lemma}\label{lem:generic}
Suppose an algorithm of format $\langle p,q,s\rangle$ and arbitrary
length is given over an infinite field, with all $U_i,V_i,W_i$ nonzero,
where $p,s\geqs1$ and $q\geqs2$. Then one can choose a nonzero column
$\eta\in K^{q\times1}$, the hyperplane
$F=\{x:x\eta=0\}\subset Q$, and nonzero rescalings of the coefficients
such that the spaces
\[
 M=\{\alpha(\eta u):u\in P\},\quad
 B_F=\{\beta(B):\im B\subseteq F\},\quad N=\im\gamma
\]
have the following properties:
\begin{enumerate}
\item $\one\in M$, and $M,B_F,N$ have full support;
\item if $a\in M$, $b\in\im\beta$, and $ab=0$, then $a=0$ or $b=0$;
\item for every $G\subseteq I$ and $i\in G$, if
$U_i\eta\in\col(\Span\{U_j:j\in G\})$, then every column of $U_i$
belongs to this same subspace.
\end{enumerate}
Here full support means that no coordinate vanishes identically on
the entire space.
\end{lemma}
\begin{proof}
Property 3 is obtained by excluding finitely many proper linear
subspaces. For fixed $G,i$, take the inverse image of
$\col(\Span U_G)$ under $\eta\mapsto U_i\eta$. If this inverse image
is proper, choose $\eta$ outside it.

Also exclude $\ker U_i$ and those $\eta$ for which
$F\subseteq\ker V_i$. The latter condition can occur only when
$\rk V_i=1$ and excludes one line of possible $\eta$.
This ensures full support for $M,B_F$. Full support for $N$ follows
from $W_i\ne0$.

For every $J$ with $\Span V_J\ne\Mat_{q,s}(K)$, choose a nonzero
matrix $B_J$ trace-orthogonal to all $V_i$, $i\in J$, and require
in addition that $B_J\eta\ne0$.
Let $u\ne0$ and $J=\supp\alpha(\eta u)$. If $\Span V_J$ were proper,
the algorithm identity would give $B_J\eta u=0$, a contradiction.
Thus $V_J$ spans the whole matrix space. Now if
$\alpha(\eta u)\beta(B)=0$, then $\beta(B)$ vanishes on $J$, so $B=0$.
This proves property 2.

Over an infinite field, finitely many proper linear subspaces cannot
cover the whole space: a product of nonzero linear equations vanishing
on them is a nonzero polynomial. Thus a suitable $\eta$ exists.
Next choose $u_0$ with $c_i=\tr(U_i\eta u_0)\ne0$ for every $i$,
and replace $U_i$ by $c_i^{-1}U_i$ and $W_i$ by $c_iW_i$.
The algorithm and all ranks are preserved, while
$\alpha(\eta u_0)=\one$. All the stated properties survive these
rescalings.
\end{proof}

An algorithm equipped with $\eta,u_0$ satisfying this lemma will be
called \emph{normalized}. Write $m(u)=\alpha(\eta u)$ and set
\[
 D=(N:M),\qquad E=MD,\qquad Z=(B_F:M),\qquad
 d=\dim D,\quad e=\dim E,\quad z=\dim Z.
\]
We call $g=e-d$ the \emph{dimension increment}.
It is nonnegative, since $\one\in M$ implies $D\subseteq E$.
For $p=q=s=3$, we have
\[
 \dim M=3,\qquad\dim B_F=6,\qquad\dim N=9,\qquad MB_F\perp N.
\]
The last orthogonality follows from $B\eta=0$ and \eqref{eq:trace}.
Lemma~\ref{lem:two-conductors} immediately gives
\begin{equation}\label{eq:dz}
 d+z\geqs45-2r.
\end{equation}

\subsection{Matrix representatives of the conductors}

\begin{definition}
For $d_0\in D$, $z_0\in Z$, and $u\in P$, define the matrices
$C_{d_0}(u)$ and $B_{z_0}(u)$ uniquely by
\[
 \gamma(C_{d_0}(u))=m(u)d_0,\qquad
 \beta(B_{z_0}(u))=m(u)z_0,\qquad \im B_{z_0}(u)\subseteq F.
\]
Existence follows from the definition of the conductor, and uniqueness
from injectivity of the coefficient maps. The dependence on each
argument is linear.
\end{definition}
\[
\begin{tikzcd}[column sep=huge,row sep=large]
 P\arrow[r,"u\mapsto C_{d_0}(u)"]\arrow[d,"m"']&
 \Mat_{p,s}(K)\arrow[d,"\gamma"]\\
 M\arrow[r,"a\mapsto ad_0"']&N
\end{tikzcd}
\qquad
\begin{tikzcd}[column sep=huge,row sep=large]
 P\arrow[r,"u\mapsto B_{z_0}(u)"]\arrow[d,"m"']&
 \{B:\im B\subseteq F\}\arrow[d,"\beta"]\\
 M\arrow[r,"a\mapsto az_0"']&B_F
\end{tikzcd}
\]

\begin{lemma}\label{lem:reciprocity}
For all $u,v\in P$, we have
\begin{align}
 vC_{d_0}(u)&=uC_{d_0}(v),\label{eq:reciprocity}\\
 C_{d_0}(u)B_{z_0}(v)&=C_{d_0}(v)B_{z_0}(u).\label{eq:compatibility}
\end{align}
Let $\E=\gamma^{-1}(E)$. For every $C\in\E$,
\begin{equation}\label{eq:symmetrizer}
 uCB_{z_0}(v)=vCB_{z_0}(u).
\end{equation}
Moreover, $B_{z_0}(u)=0$ implies $z_0=0$ or $u=0$.
\end{lemma}
\begin{proof}
Substitute $A=\eta v$ and $\gamma(C)=m(u)d_0$ in \eqref{eq:trace}.
The right-hand side is symmetric in $u,v$; varying $B\eta$ gives
\eqref{eq:reciprocity}. Similarly, the product $m(u)m(v)d_0z_0$
and trace duality give \eqref{eq:compatibility}.
For $C=C_{d_0}(w)$, we obtain
\[
 uC_{d_0}(w)B_{z_0}(v)
 =wC_{d_0}(u)B_{z_0}(v)
 =wC_{d_0}(v)B_{z_0}(u)
 =vC_{d_0}(w)B_{z_0}(u).
\]
These $C$ span $\E$, proving \eqref{eq:symmetrizer}.
The last assertion follows from property 2 of Lemma~\ref{lem:generic}.
\end{proof}

\begin{lemma}\label{lem:physical-bound}
For a normalized algorithm of format $\langle p,q,s\rangle$ and
arbitrary length $r$,
\[
 ps+(q-1)e\leqs r.
\]
In particular, if $p=q=s=3$ and $r\leqs20$, then $e\leqs5$.
\end{lemma}
\begin{proof}
Consider the $r$ linear maps $Q_i:T\mapsto V_iTU_i\eta$ from
$\Mat_{s,p}(K)$ to the column space $K^q$.
Identities \eqref{eq:sandwich} and \eqref{eq:reciprocity} imply
\begin{align*}
 \sum_i\gamma(C)_iQ_i(T)&=\tr(TC)\eta,\\
 \sum_i(d_0)_i\alpha(\theta u)_iQ_i(T)
       &=\tr(TC_{d_0}(u))\theta.
\end{align*}
For the second identity it suffices to take $T=av$ of rank one.
By \eqref{eq:trace}, its left-hand side is
$(vC_{d_0}(u)a)\theta$, as required; these matrices span the whole space.
Thus the span of the $Q_i$ contains every map $T\mapsto\tr(TC)\eta$
and every map $T\mapsto\tr(TC)\theta$ with $C\in\E$, $\theta\in K^q$.
These two spaces have dimensions $ps$ and $qe$, and their intersection
has dimension $e$. Their sum therefore has dimension
$ps+(q-1)e\leqs r$.
\end{proof}

\Needspace{10\baselineskip}
\begin{cor}\label{cor:two-two-n}
For every field $K$ and integer $n\geqs1$,
\[
 R_K\langle2,2,n\rangle\geqs\left\lceil\frac{7n}{2}\right\rceil.
\]
\end{cor}
\begin{proof}
An algorithm over $K$ remains valid over an algebraic closure, and
deleting zero summands does not increase its length. It therefore
suffices to bound the length $r$ of an algorithm with no zero summands
over an infinite field. The restriction $r\leqs20$ is not assumed here.

Apply Lemma~\ref{lem:generic} with $p=q=2$ and $s=n$.
Injectivity of the coefficient maps gives
\[
 \dim M=2,\qquad \dim N=2n,\qquad \dim B_F=n.
\]
Since $\one\in M$, choose $x$ so that $M=\Span(\one,x)$.
The relation $MB_F\perp N$ and commutativity of the coordinatewise
product imply $MN\perp B_F$. Nondegeneracy of the pairing gives
\[
 \dim MN\leqs r-n.
\]
Consider the surjective linear map
\[
 \Phi:N\oplus N\longrightarrow MN=N+xN,\qquad
 (u,v)\longmapsto xu-v.
\]
Its kernel consists of the pairs $(u,xu)$ with $u\in N$ and $xu\in N$.
This is exactly the condition $u\in D=(N:M)$. By rank-nullity,
\[
 d=\dim D=4n-\dim MN\geqs5n-r.
\]
Since $\one\in M$, we have $D\subseteq MD$, hence $e\geqs d$.
Lemma~\ref{lem:physical-bound} now gives
\[
 r\geqs2n+e\geqs2n+d\geqs7n-r.
\]
Thus $2r\geqs7n$, and the integrality of $r$ completes the proof.
\end{proof}

\subsection{Joint symmetrizers}

\begin{definition}
Suppose $B_z(u):S\to F$ depends bilinearly on $z\in Z$ and $u\in P$.
Its \emph{space of joint symmetrizers} is
\[
 \Sig(B)=\{C:P\to S:\ uCB_z(v)=vCB_z(u)
                          \text{ for all }u,v,z\}.
\]
In chosen bases, this is the solution space of linear conditions
requiring certain matrix products to be symmetric.
Lemma~\ref{lem:reciprocity} gives $\E\subseteq\Sig(B)$.
\end{definition}

\begin{lemma}\label{lem:rank-one-space}
If every matrix in a linear space has rank at most one, then the space
has either a common one-dimensional column support or a common
one-dimensional row support.
\end{lemma}
\begin{proof}
Choose a nonzero matrix $ab$. For another nonzero matrix $cd$, the sum
$ab+cd$ has rank two if both $a,c$ and $b,d$ are independent.
Thus every element has its first factor proportional to $a$ or its
second factor proportional to $b$. If both types occur with
nonproportional free factors, their sum again has rank two.
Hence one factor is common to the entire space.
\end{proof}

\begin{lemma}\label{lem:sym-small}
Let $\dim P=\dim S=3$ and $\dim F=2$, and suppose that $B_z(u)=0$
implies $z=0$ or $u=0$. If $\dim Z\geqs3$ and
$\dim\Span\{B_z(u):z,u\}\geqs5$, then $\dim\Sig(B)\leqs1$.
\end{lemma}
\begin{proof}
Let $C\in\Sig(B)$ be singular, and choose $0\ne u\in\ker C$.
Symmetry implies that all $B_z(u)$ annihilate $\im C$.
The map $z\mapsto B_z(u)$ is injective, so
\begin{equation}\label{eq:sym-capacity}
 \dim Z\leqs(3-\rk C)\dim F.
\end{equation}
Consequently, a singular symmetrizer has rank at most one.

Suppose there are an invertible symmetrizer $C_0$ and a nonzero
rank-one symmetrizer $C$. Choose a row $v$ with $vC\ne0$ and set
$x=(vC)C_0^{-1}$. For $u\in\ker C$, the two symmetry identities give
\[
 uC_0B_z(x)=xC_0B_z(u)=vCB_z(u)=uCB_z(v)=0.
\]
Thus every $B_z(x)$ annihilates the two-dimensional space $(\ker C)C_0$.
Injectivity in $z$ gives $\dim Z\leqs2$, a contradiction.
Therefore, if $C_0$ exists, every nonzero element of $\Sig(B)$ is
invertible. There cannot be two independent such elements: the matrix
$C_0^{-1}C_1$ has an eigenvalue, making $C_1-\lambda C_0$ a nonzero
singular element.

It remains to consider the case where all symmetrizers have rank at
most one. Suppose their space $H$ has dimension at least two.
Lemma~\ref{lem:rank-one-space} gives two possibilities.
If all $C\in H$ have image in a common line $L\subset S$, write two
independent elements as $a_1s,a_2s$, with $Ks=L$.
Symmetry says that $u\mapsto sB_z(u)$ is proportional to each of the
independent forms $u\mapsto ua_1$ and $u\mapsto ua_2$.
It must therefore vanish: $LB_z(u)=0$ for all $z,u$.
The span of all $B_z(u)$ then has dimension at most
$\dim\Hom(S/L,F)=4$.
If instead all $C\in H$ have a common kernel hyperplane $H_0\subset P$,
then for $x\notin H_0$ the map $C\mapsto xC$ is injective.
Its image has dimension $\dim H\geqs2$. For a nonzero $u\in H_0$,
symmetry forces all $B_z(u)$ to annihilate this image, again giving
$\dim Z\leqs(3-\dim H)2\leqs2$. Both possibilities are impossible.
\end{proof}

\begin{lemma}\label{lem:five-sym}
Under the assumptions of the previous lemma, except for the bounds on
$\dim Z$ and the span of $B$, suppose that $Z\ne0$ and
$\dim\Sig(B)\geqs5$. Then there are an invertible matrix $L$ and an
injective linear map $f:Z\to F$ such that
\begin{equation}\label{eq:factor}
 B_z(u)= (Lu^{\mathsf t}) f(z).
\end{equation}
In other words, its value on a row $s$ is $(sLu^{\mathsf t})f(z)$.
\end{lemma}
\begin{proof}
First establish an auxiliary fact. If independent matrices
$A,B\in\Mat_3(K)$ satisfy $Ak=Bk=0\Rightarrow k=0$ for columns $k$,
then the space of matrices $T$ for which both $TA$ and $TB$ are
symmetric has dimension at most four. The dual equations are
represented by the map
\[
 (k,l)\longmapsto A[k]_\times+B[l]_\times,\qquad
 [k]_\times=\begin{pmatrix}0&-k_3&k_2\\k_3&0&-k_1\\-k_2&k_1&0\end{pmatrix}.
\]
Here $k,l$ are columns, and we write $k\times l=[k]_\times l$ for this
formula over $K$. If the image is zero, multiplication on the right by
$k$ and $l$ shows that $A(k\times l)=B(k\times l)=0$.
Hence $k,l$ are dependent. Regard the kernel as a space of
$3\times2$ matrices $(k\ l)$: every matrix in it has rank at most one.
If its dimension is at least two, Lemma~\ref{lem:rank-one-space} gives
a common first or second factor. In the first case, the kernel contains
$(k,0),(0,k)$ with $k\ne0$, so $A[k]_\times=B[k]_\times=0$, contrary
to the trivial common kernel. In the second case it consists of
$(ak,bk)$ with fixed $(a,b)\ne(0,0)$. The nonzero matrix $aA+bB$
can annihilate $[k]_\times$ only for $k$ on a single line: any nonzero
row of that matrix must be proportional to $k$.
Thus the kernel has dimension at most one, and the map has rank at
least five. The required space of symmetrizers has codimension at
least five. This argument uses the zero diagonal of $[k]_\times$
and does not divide by two.

Choose $z\ne0$ and represent the two output coordinates of the
bilinear map $(u,s)\mapsto sB_z(u)$ by $3\times3$ matrices $A,B$.
The expression is $(sAu^{\mathsf t},sBu^{\mathsf t})$.
Since $u\mapsto B_z(u)$ is injective, $Ak=Bk=0$ implies $k=0$.
The existence of five symmetrizers and the auxiliary fact force $A,B$
to be proportional. Their common nonzero matrix factor $L$ is
invertible. Compare it with each coordinate matrix for the other
values of $z$: the common kernel is again trivial, and the same fact
forces proportionality to $L$. The proportionality coefficients
define a linear map $f$. The absence of nonzero pairs with zero
product makes $f$ injective.
\end{proof}

\begin{lemma}\label{lem:injective-product}
If the family representing the second conductor has form
\eqref{eq:factor}, then $\dim NZ=9z$ and $r\geqs9z+3$.
The same conclusion holds if there are row bases $(u_j)$ of $P$ and
$(s_j)$ of $S$, and injective maps $f_j:Z\to F$, such that
\begin{equation}\label{eq:diagonal-factor}
 B_z(u_j)=s_j^*f_j(z).
\end{equation}
Here $s_j^*$ is the corresponding dual-basis column:
$s_k s_j^*=\delta_{kj}$.
\end{lemma}
\begin{proof}
Let $z_1,\ldots,z_z$ be a basis of $Z$, and suppose
$\sum_a\gamma(C_a)z_a=0$. Multiply by $m(u)$, use the definition of
$B_{z_a}(u)$, and then apply \eqref{eq:cycles}. This gives
\[
 \sum_a C_a B_{z_a}(u)=0\qquad\text{for every }u.
\]
In case \eqref{eq:factor}, this is $\sum_a C_a Lu^{\mathsf t}f(z_a)=0$.
Independence of the rows $f(z_a)$ gives $C_aLu^{\mathsf t}=0$ for
all $a,u$, hence $C_a=0$. In case \eqref{eq:diagonal-factor}, likewise
$C_a s_j^*=0$ for all $a,j$, again giving $C_a=0$.
Thus the map $(C_a)_a\mapsto\sum_a\gamma(C_a)z_a$ is injective.
Its image is $NZ$, which consequently has dimension $9z$.
Since $Z\subseteq B_F$ and $MB_F\perp N$, we have $M\perp NZ$.
Therefore $r\geqs\dim M+\dim NZ=3+9z$.
\end{proof}

\begin{prop}\label{prop:critical}
For a normalized algorithm of format $3\times3\times3$ and length
$r\leqs20$, we have
\[
 3\leqs d\leqs e\leqs5,\qquad z\leqs2,\qquad d+z\geqs45-2r,
 \qquad e=5\Longrightarrow d\geqs4, z\leqs1.
\]
In particular, $r\geqs20$; when $r=20$, the increment $g=e-d$ is zero
or one.
\end{prop}
\begin{proof}
Lemma~\ref{lem:nz-growth}, applied to $M,Z$, gives
$\dim MZ\geqs z+2$ if $Z\ne0$. Since $MZ\subseteq B_F$, we get
$z\leqs4$. Together with \eqref{eq:dz}, this gives $d\geqs1$.
By Lemma~\ref{lem:physical-bound}, $e\leqs5$.

If $z\geqs3$, the span of the matrices $B_z(u)$ has dimension
$\dim MZ\geqs5$. Lemma~\ref{lem:sym-small} gives $e\leqs1$, so
$D=E$ is one-dimensional. Its nonzero vector $w=\gamma(C)$ is a common
eigenvector for multiplication by $M$. Choose an active coordinate $i$;
the space $M_0=\{m\in M:m_i=0\}$ has dimension two.
For $m(u)\in M_0$, we have $C_w(u)=0$.
Using \eqref{eq:compatibility} with $u_0$ satisfying $m(u_0)=\one$
gives $CB_z(u)=0$. Thus the matrix representatives of $M_0Z$ lie in
$\Hom(S/\im C,F)$, whose dimension is at most four.
Lemma~\ref{lem:nz-growth} gives $z+1\leqs4$, hence $z\leqs3$.
Now $d+z\leqs4$, contradicting \eqref{eq:dz}. Therefore $z\leqs2$
and $d\geqs3$.

If $e=5$ and $z=2$, Lemmas~\ref{lem:five-sym} and
\ref{lem:injective-product} give $r\geqs21$, a contradiction.
Hence $z\leqs1$ and $d\geqs4$ in this case.
If $r\leqs19$, then \eqref{eq:dz} would give $d+z\geqs7$;
this is impossible when $e\leqs4$, while $e=5$ already gives
$d+z\leqs6$. Thus $r=20$. Every remaining pair $(d,e)$ has difference
zero or one.
\end{proof}

\section{Supports of tests and atomic tests}

\subsection{Two bounds for tests with prescribed support}

\begin{definition}
For $G\subseteq I$, set
\[
 \test_G=\{C\in\Mat_{p,s}(K):\supp\gamma(C)\subseteq G\},\quad
 \delta_G=\dim\test_G,\quad
 H_G=\sum_{C\in\test_G}\im C,\quad h_G=\dim H_G.
\]
Thus $\test_G$ consists of the tests supported in $G$, and $H_G$ is
their common row support.
\end{definition}

\begin{lemma}\label{lem:budgets}
Suppose the columns of all $U_i$, $i\in G$, lie in a fixed
$k$-dimensional space $H\subset K^{p\times1}$, with $k<p$, and let
$m=|G|$. Then the columns of every $C\in\test_G$ also lie in $H$, and
\begin{align}
 \delta_G&\leqs k h_G,\label{eq:capacity}\\
 \delta_G+(q-1)h_G&\leqs m,\label{eq:internal}\\
 m+qs+ks&\leqs r+\delta_G.\label{eq:external}
\end{align}
Inequality \eqref{eq:internal} requires the normalization
$\alpha(\eta u_0)=\one$. In addition,
\begin{equation}\label{eq:mass}
 \Span\{V_i:i\in G\}\supseteq\{V: \im V\subseteq H_G\},
 \qquad \dim\Span V_G\geqs qh_G.
\end{equation}
\end{lemma}
\begin{proof}
For $C\in\test_G$, identities \eqref{eq:cycles} give
$CB\in\Span U_G$ and $AC\in\Span V_G$ for all $A,B$.
The first inclusion puts the columns of $C$ in $H$; the second gives
\eqref{eq:mass}. Inequality \eqref{eq:capacity} follows from the
dimension of the space of matrices with columns in $H$ and rows in $H_G$.

Let $f:K^G\to\Mat_{q,s}(K)$ be given by $f(c)=\sum_{i\in G}c_iV_i$.
The injective map $\gamma:\test_G\to K^G$ satisfies
$f\gamma(C)=\eta u_0C$, so $\dim f(\gamma\test_G)\leqs h_G$.
Combining
\[
 \dim\im f+\dim\gamma\test_G\leqs m+\dim f(\gamma\test_G)
\]
with \eqref{eq:mass} gives \eqref{eq:internal}.

For \eqref{eq:external}, restrict the first input by $AH=0$.
All summands indexed by $G$ disappear. On the space of tests with
columns in $H$, restricting $\gamma$ to $I\setminus G$ has kernel
$\test_G$, hence rank $ks-\delta_G$. Choose that many independent
coordinate forms. By adding to any test a matrix with columns in $H$,
we can make these forms zero; the product in \eqref{eq:trace} is
unchanged because $AH=0$. We obtain an algorithm for rectangular
matrix multiplication with inner dimension $p-k>0$ and at most
$r-m-ks+\delta_G$ summands. Its coefficient map on the space of
$q\times s$ matrices is injective, so its length is at least $qs$.
This is \eqref{eq:external}.
\end{proof}

\begin{lemma}\label{lem:rect15}
Bilinear matrix multiplication of format $\langle3,2,3\rangle$ requires
at least fifteen summands.
\end{lemma}
\begin{proof}
Apply the normalization and conductor construction with $p=s=3$,
$q=2$. Then $\dim M=3$, $\dim B_F=3$, and $\dim N=9$.
If $Z\ne0$, Lemma~\ref{lem:nz-growth} gives $z\leqs1$; if $Z=0$,
the same bound is immediate. Consequently,
$d\geqs36-2r-z\geqs35-2r$.
Lemma~\ref{lem:physical-bound} gives $9+e\leqs r$, and $d\leqs e$.
Thus $44\leqs3r$, hence $r\geqs15$.
Passing to an infinite extension field and deleting zero summands do
not increase the length.
\end{proof}

\begin{lemma}\label{lem:point-capacity}
In an algorithm of format $\langle3,3,3\rangle$ and length twenty,
at most three matrices $U_i$ have their columns in any fixed line.
If there are exactly three such indices, then $\delta_G=1$, and their
$U_i$ span the full three-dimensional space of matrices with columns
in that line. The same holds for rows and for cyclic permutations
of $U,V,W$.
\end{lemma}
\begin{proof}
Here $k=1$. Tests supported in $G$ have a fixed first factor, so
$\delta_G=h_G$, and \eqref{eq:mass} gives $3\delta_G\leqs m$.
In the proof of \eqref{eq:external}, deleting $G$ and a further
$3-\delta_G$ terms leaves format $\langle3,2,3\rangle$.
Lemma~\ref{lem:rect15} gives
\[
 15\leqs20-m-3+\delta_G,\qquad m\leqs\delta_G+2.
\]
Together these bounds yield $\delta_G\leqs1$ and $m\leqs3$.
When $m=3$, we get $\delta_G=1$.
For nonzero $C=as\in\test_G$, the matrices $CB$, with $B$ arbitrary,
span the whole space $\{ab:b\in K^{1\times3}\}$ and lie in the span
of $U_G$ by \eqref{eq:cycles}. This proves the assertion about the span.
Transposition and cyclic permutation in \eqref{eq:trace} give the
other versions.
\end{proof}

\begin{lemma}\label{lem:partial}
Suppose a linear space of output forms of dimension $\delta\geqs3$
is computed from the product of $2\times n$ and $n\times2$ matrices.
Then the number of summands is at least $3n+\delta-3$.
\end{lemma}
\begin{proof}
First let $\delta=3$. A three-dimensional space of $2\times2$ test
matrices is the trace annihilator of one nonzero matrix. Independent
changes of the two bases reduce this matrix to one of two rank cases.
In the rank-one case the three computed forms are
$x u^{\mathsf t},x v^{\mathsf t},y u^{\mathsf t}$, where $x,y,u,v\in K^n$.
The first-input forms evaluated on $(0,y)$ separate all $y$; otherwise
the third output form would vanish identically for some nonzero $y$.
Choose $n$ independent such forms and make them zero by a linear
substitution $y=Lx$. At most $r-n$ terms remain, computing the first
two outputs. They separate all $2n$ coordinates of $(u,v)$, so
$r-n\geqs2n$.

In the rank-two case the outputs are
$x u^{\mathsf t}-y v^{\mathsf t},x v^{\mathsf t},y u^{\mathsf t}$.
Their matrices on $K^n\oplus K^n$ are
\[
 S_0=\begin{pmatrix}I&0\\0&-I\end{pmatrix},\quad
 S_1=\begin{pmatrix}0&I\\0&0\end{pmatrix},\quad
 S_2=\begin{pmatrix}0&0\\I&0\end{pmatrix}.
\]
The block map
\[
 (v_0,v_1,v_2)\longmapsto
 (S_1v_2-S_2v_1,\ S_2v_0-S_0v_2,\ S_0v_1-S_1v_0)
\]
has trivial kernel: use the last two equations to express $v_1,v_2$
in terms of $v_0$, and the first becomes
$(S_1S_0^{-1}S_2-S_2S_0^{-1}S_1)v_0=0$; the operator in parentheses
is $-I$. The rank is therefore $6n$.
A single bilinear summand contributes the same rank-one matrix to
all three $S_a$, with three scalar coefficients. Its contribution
to the block map thus has rank at most that of the corresponding
alternating $3\times3$ matrix, which is at most two.
Consequently, $2r\geqs6n$. The sign causes no problem in characteristic two.

If $\delta>3$, the coefficient forms embed the test space injectively
into $K^r$. Choose $\delta-3$ independent coordinate forms and take
their common kernel, of dimension three. On it, $\delta-3$ summands
disappear, and the bound just proved gives
$r-(\delta-3)\geqs3n$.
\end{proof}

\begin{lemma}\label{lem:plane}
In an algorithm of length twenty, at most seven $U_i$ have their
columns in any fixed plane $H\subset K^3$. For every set of such
indices, $\delta_G\leqs2$.
\end{lemma}
\begin{proof}
First normalize the algorithm using Lemma~\ref{lem:generic}; nonzero
rescalings preserve the supports and dimensions in question.
Lemma~\ref{lem:budgets} with $k=2$ gives
$m\leqs\delta_G+5$, $\delta_G+2h_G\leqs m$, and
$\delta_G\leqs2h_G$. Hence $h_G\leqs2$.
If $\delta_G\geqs3$, then $h_G=2$. The tests in $\test_G$ restrict
the product of $2\times3$ and $3\times2$ matrices and compute
$\delta_G$ independent output forms. Lemma~\ref{lem:partial} gives
$m\geqs\delta_G+6$, a contradiction. Thus $\delta_G\leqs2$ and $m\leqs7$.
\end{proof}

\subsection{Atomic tests}

\begin{definition}
An \emph{atomic test} is a nonzero rank-one matrix $C=as$, where $a$
is a column and $s$ a row, such that $|\supp\gamma(C)|=3$.
The set $G=\supp\gamma(C)$ is its support; $Ka$ and $Ks$ are its
column and row directions.
\end{definition}

\begin{lemma}\label{lem:atom}
For every nonzero test $C$, we have $|\supp\gamma(C)|\geqs3\rk C$.
For an atomic test $C=as$ with support $G$, there are rows $b_i$ and
columns $c_i$ such that
\[
 U_i=ab_i,\qquad V_i=c_i s,
 \qquad \gamma(C)_i\, b_jc_i=\delta_{ji}\quad(i,j\in G).
\]
In particular, every $U_i,V_i$ on $G$ has rank one.
The supports of two atomic tests are either equal or disjoint.
In the first case the tests are proportional. The column directions
of any three tests with distinct supports are independent, as are
their row directions.
\end{lemma}
\begin{proof}
The bilinear form $(A,B)\mapsto\tr(CBA)$ has matrix rank $3\rk C$:
after reducing $C$ to rank normal form, it is a sum of $3\rk C$
independent coordinate products. Representation \eqref{eq:trace}
contains only $|\supp\gamma(C)|$ rank-one forms, proving the bound.

For $C=as$ and exactly three terms, equality holds in the rank bound.
Factoring a full-rank bilinear form through $K^3$ shows that the
first coefficient forms constitute a basis of the forms $A\mapsto bAa$,
and the second a basis of the forms $B\mapsto sBc$.
Hence $U_i=ab_i$ and $V_i=c_i s$. Equality with the matrix pairing
gives the stated weighted duality.

If two atoms share an index, its nonzero $U_i$ fixes their common
column direction. Lemma~\ref{lem:point-capacity} permits only three
indices in this class, so the supports coincide. The space of tests
with this support is one-dimensional, so the tests are proportional.
Finally, dependence of the column directions of three distinct atoms
would put their nine $U_i$ in one plane, contrary to Lemma~\ref{lem:plane}.
For row directions, use the symmetry involving transposition.
\end{proof}

\begin{lemma}\label{lem:stationary}
Let $0\ne d_0\in D$, $G=\supp d_0$, and $a=\dim(Md_0)$.
Then there is an $a$-dimensional space of columns $H$ such that
\[
 \Span U_G=\{T:\text{all columns of }T\text{ lie in }H\}.
\]
In particular, $\dim\Span U_G=3a$.
If $w\in N\setminus\{0\}$ is a common eigenvector for multiplication
by $M$, then $w=\gamma(C)$ for an atomic test $C$.
\end{lemma}
\begin{proof}
Let $H$ be the span of the columns of all $C_{d_0}(u)$.
Reciprocity \eqref{eq:reciprocity} gives
\[
 H=(\ker(u\mapsto C_{d_0}(u)))^\perp:
\]
a row $v$ annihilates all
these columns precisely when $uC_{d_0}(v)=0$ for every $u$.
Thus $\dim H=\dim Md_0=a$.

All $C_{d_0}(u)$ are supported in $G$, so \eqref{eq:cycles} implies
$H\subseteq\col(\Span U_G)$. The equality
$\gamma(C_{d_0}(u))_i=(uU_i\eta)(d_0)_i$ shows that $U_i\eta\in H$
for $i\in G$: both sides vanish on the kernel of the family $C_{d_0}$.
Property 3 of Lemma~\ref{lem:generic} now puts every column of $U_i$
in $\col(\Span U_G)$. This latter space is itself contained in $H$,
since its elements can be obtained by evaluating matrices in
$\Span U_G$ on $\eta\ne0$. Therefore every column of $U_i$ lies in $H$;
the reverse inclusion of the full matrix space is already part of
the definition of $\col$.

For a common eigenvector $w$, we have $w\in D$ and $\dim Mw=1$.
Thus all $U_i$ on its support have one column direction.
By Lemma~\ref{lem:point-capacity} the support has at most three
indices, and by Lemma~\ref{lem:atom} it has at least three.
Its test has rank one and is atomic.
\end{proof}

\section{Excluding a positive dimension increment}

\subsection{Local geometry of a growing component}

\begin{lemma}\label{lem:growing-block}
If $e=d+1$, then $E$ has a coordinate component $G$ with the following
properties:
\begin{enumerate}
\item $\dim(M|_G)=2$, $\dim D[G]=1$,
$E[G]=N[G]$, and $\dim E[G]=2$;
\item $6\leqs |G|\leqs7$, and the span of $U_G$ is the full space of
matrices whose columns lie in a certain plane;
\item $\test_G$ is two-dimensional, its column and row supports are
two-dimensional, and it contains a rank-two matrix;
\item in $G$, at most two matrices $U_i$ have any fixed column
direction, and at most two have any fixed row direction;
\item $N[G]$ contains no common eigenvector for multiplication by $M$.
\end{enumerate}
\end{lemma}
\begin{proof}
Lemma~\ref{lem:component-cost} gives a component on which the restriction
of $M$ has dimension two. There is a vector $d_0\in D[G]$ with support
exactly $G$: the supports of $D[G]$ and $E[G]$ coincide, and over an
infinite field we can avoid every coordinate hyperplane.
By Lemma~\ref{lem:stationary}, the span of $U_G$ is the full space of
matrices whose columns lie in a certain plane. Its dimension is six,
so $|G|\geqs6$. Lemma~\ref{lem:plane} gives $|G|\leqs7$ and
$\dim N[G]\leqs2$. Now
\[
 2+\dim D[G]\leqs\dim E[G]+1,\qquad
 1\leqs\dim D[G],\qquad E[G]\subseteq N[G]
\]
force the dimensions to be $1,2,2$. If $N[G]$ contained a common
eigenvector $w$, it would lie in $D[G]$ and span it; then
$MD[G]=Kw$, contrary to $\dim E[G]=2$.

The column support of $\test_G$ is two-dimensional by the proof of
Lemma~\ref{lem:stationary}. Inequality \eqref{eq:internal} gives
$2+2h_G\leqs7$, hence $h_G\leqs2$.
We show that $h_G\ne1$. If every test has its rows in a line $L$,
then $\test_G=\{as:a\in H\}$ with $\dim H=2$ and $Ks=L$.
Consider the maps
\[
 f_i:\test_G\to\Mat_{3,3}(K),\qquad f_i(C)=\gamma(C)_iV_i.
\]
Their span contains all maps $C\mapsto AC$, which form a space of
dimension $3\dim H=6$. Projecting the values modulo the space of
matrices whose rows lie in $L$ kills these six dimensions.
Thus the span of the maps $C\mapsto\gamma(C)_i[V_i]$ has dimension
at most $|G|-6\leqs1$.
For indices with $[V_i]\ne0$, all the nonzero forms
$\gamma(\cdot)_i|_{\test_G}$ are consequently proportional: evaluate
two proportional maps on any test where one of them is nonzero.
There is a nonzero $C\in\test_G$ annihilating all these forms.
Its support consists only of indices for which the rows of $V_i$
lie in $L$. By Lemma~\ref{lem:point-capacity} there are at most three
such indices, so $C$ is atomic.
For an atom, $U_i=ab_i$; normalization gives $(u_0a)b_i\eta=1$ on
its support. Hence $m(u)_i=(ua)/(u_0a)$ is constant there.
Thus $\gamma(C)$ is a common eigenvector, which has already been excluded.

Both supports are now two-dimensional. If every test had rank at
most one, Lemma~\ref{lem:rank-one-space} would make one support
one-dimensional. Thus a rank-two test exists.

Three indices with a common column direction would support a nonzero
test by Lemma~\ref{lem:point-capacity}. This would be an atom and a
common eigenvector, a contradiction. For a row direction, apply the
same lemma after transposing: three matrices with a common row
direction must span the three-dimensional space of all matrices
with that direction. But the additional requirement that their
columns lie in a plane leaves only two dimensions, a contradiction.
\end{proof}

\subsection{Replacing one bilinear pair}

\begin{definition}
The \emph{first rank sum} of an algorithm is $\sigma_U=\sum_i\rk U_i$.
Among all algorithms of length twenty over the chosen field, choose
one minimizing $\sigma_U$. A minimum exists because the possible
values form a nonempty set of natural numbers. Normalization by
nonzero rescalings preserves it.
The \emph{bilinear pair} $U\boxtimes V$ denotes the form
$(A,B)\mapsto\tr(UA)\tr(VB)$, not the matrix product $UV$.
\end{definition}

\begin{lemma}\label{lem:exchange}
In an algorithm with minimal $\sigma_U$, an equality
\[
 \sum_i a_i(U_i\boxtimes V_i)=X\boxtimes Y,\qquad a_j\ne0,
\]
implies $\rk U_j\leqs\rk X$.
\end{lemma}
\begin{proof}
Solve for $U_j\boxtimes V_j$ and substitute into the algorithm.
The $j$th term is replaced by the pair $(X,Y)$ with output
$a_j^{-1}W_j$; for $i\ne j$, the output becomes
$W_i-a_i a_j^{-1}W_j$. All other first coefficient matrices are
unchanged, and the length remains twenty. If $\rk X<\rk U_j$,
then $\sigma_U$ decreases.
\end{proof}

\begin{prop}\label{prop:no-growth}
For an algorithm with minimal first rank sum, the normalization in
Lemma~\ref{lem:generic} satisfies $MD=D$.
\end{prop}
\begin{proof}
By Proposition~\ref{prop:critical}, it suffices to exclude $e=d+1$.
Take the component $G$ from Lemma~\ref{lem:growing-block}.
Compress its two outer supports to $K^2$. The first coefficients have
form $U_i=JX_i$, where $J$ is a $3\times2$ matrix of rank two and
$X_i$ are $2\times3$ matrices. Choose a basis $C,D$ of the test space
such that, on these two supports, $C$ is the identity matrix and
$D$ is a nonscalar $2\times2$ matrix $H$. We may also require
\[
 c_i=\gamma(C)_i\ne0\quad(i\in G):
\]
each coordinate on $\test_G$ is nonzero, while invertibility and
these finitely many conditions are expressed by nonzero polynomials
in the two test coordinates. Set $t_i=\gamma(D)_i/c_i$.

Subtracting the two versions of \eqref{eq:cycles} for $C,D$ gives
a matrix relation with coefficients multiplying the $V_i$:
\begin{equation}\label{eq:cross-relation}
 \sum_{i\in G} c_i(t_iI-H)X_i\boxtimes V_i=0.
\end{equation}
This equality is interpreted coordinatewise in the matrix of the
first factor. It also follows directly from
$\eta'(uC)=\sum_i\gamma(C)_i(uU_i\eta')V_i$ for arbitrary $u,\eta'$.

By \eqref{eq:mass}, $\dim\Span V_G\geqs6$.
If $V_G$ were independent, every first coefficient in
\eqref{eq:cross-relation} would vanish. Then the columns of each
$X_i$ would lie in an eigenline of $H$ with eigenvalue $t_i$.
A nonscalar $2\times2$ matrix over an algebraically closed field has
at most two eigendirections. By Lemma~\ref{lem:growing-block}, each
contains at most two indices, so $|G|\leqs4$, a contradiction.
Consequently,
\[
 |G|=7,\qquad \dim\Span V_G=6.
\]
All relations among the $V_i$ are spanned by one nonzero vector $w$,
with $\sum_iw_iV_i=0$. Each coordinate of the matrix relation
\eqref{eq:cross-relation} is a multiple of $w$. Thus there is a
$2\times3$ matrix $T$ satisfying
\begin{equation}\label{eq:resolvent}
 (t_iI-H)X_i=(w_i/c_i)T\qquad(i\in G).
\end{equation}

We claim that $\rk T=2$. Let $\lambda,\mu$ be the roots of the
characteristic polynomial of $H$, possibly equal. Both shifts
$H-\lambda I$ and $H-\mu I$ are nonzero of rank one, and
\begin{equation}\label{eq:CH}
 (H-\lambda I)(H-\mu I)=0.
\end{equation}
If $T=0$, the preceding count again gives $|G|\leqs4$.
Suppose $T=ab$ has rank one. If $a$ is not an eigenvector column of
$H$, then at $t_i=\lambda,\mu$, multiplying \eqref{eq:resolvent} by
the opposite shift gives $w_i=0$ and puts the columns of $X_i$ in
the corresponding eigenline. Away from these two parameters the
shift is invertible, and all $X_i$ have the common row direction $Kb$.
Each of these three classes has at most two indices, giving at most
six in total, a contradiction.

If $Ha=\lambda a$, multiplying \eqref{eq:resolvent} by $H-\lambda I$
and using \eqref{eq:CH} shows that, for $t_i\ne\mu$, all columns of
$X_i$ lie in $Ka$. But these indices are exactly the support of the
nonzero test $D-\mu C$. There are at most two of them by
Lemma~\ref{lem:growing-block}, whereas every nonzero test requires at
least three indices. The case $Ha=\mu a$ is analogous.
Therefore $\rk T=2$.

At the roots $t_i=\lambda,\mu$, the left-hand side of
\eqref{eq:resolvent} has rank at most one, so $w_i=0$.
Set $J_0=\{i\in G:t_i\ne\lambda,\ t_i\ne\mu\}$.
For $i\in J_0$, equations \eqref{eq:CH} and \eqref{eq:resolvent} give
\[
 c_i(t_i-\lambda)X_i
 =w_iT+\frac{w_i}{t_i-\mu}(H-\lambda I)T.
\]
Pair this with the second bilinear factor $V_i$ and sum.
The first term vanishes because $\sum_iw_iV_i=0$ and $w_i=0$
outside $J_0$:
\[
 \sum_{i\in J_0}c_i(t_i-\lambda)(X_i\boxtimes V_i)
 =((H-\lambda I)T)\boxtimes
     \left(\sum_{i\in J_0}\frac{w_i}{t_i-\mu}V_i\right).
\]
Multiplying the first matrices on the left by $J$ transfers the
identity to the original $U_i$. Choose $j$ with $w_j\ne0$.
Then $j\in J_0$, the coefficient of its pair is nonzero, and
$\rk U_j\geqs\rk X_j\geqs\rk T=2$ by \eqref{eq:resolvent}.
But the first factor on the right, $J(H-\lambda I)T$, has rank at
most one. This contradicts Lemma~\ref{lem:exchange}.
\end{proof}

\section{Four atoms: a residual sum of eight terms}

\subsection{A five-dimensional annihilator}

\begin{lemma}\label{lem:annihilator}
Let $C_j=a_js_j$, $1\leqs j\leqs4$, be rank-one $3\times3$ matrices,
with every three $a_j$ independent and every three $s_j$ independent.
Then the space
\[
 L=\{T\in\Mat_3(K):\tr(TC_j)=0\ (1\leqs j\leqs4)\}
\]
has dimension five and admits a basis of rank-one matrices.
\end{lemma}
\begin{proof}
Independent changes of basis in the two outer spaces and nonzero
rescalings allow the column representatives to be chosen as
\[
 e_1,\quad e_2,\quad e_3,\quad e_1+e_2+e_3.
\]
Independently, the row representatives can be reduced to rows with
the same coordinates. The annihilator is therefore defined by
\[
 T_{11}=T_{22}=T_{33}=0,\qquad \sum_{i\ne j}T_{ij}=0.
\]
These are four independent equations. Consider the six rank-one matrices
\[
 E_{12}-E_{13},\ E_{21}-E_{23},\ E_{31}-E_{32},\quad
 E_{21}-E_{31},\ E_{12}-E_{32},\ E_{13}-E_{23}.
\]
They span this hyperplane in the six-dimensional space of off-diagonal
coordinates: they are differences of adjacent coordinates along a
connected six-cycle. Any five that leave a connected path form a basis.
Reversing the changes of basis preserves the ranks.
\end{proof}

\begin{prop}\label{prop:residual-model}
If an algorithm of length twenty has four atomic tests with distinct
supports, then there exist a five-dimensional space $L$, a basis
$b_1,\ldots,b_5$ of $L$, eight vectors $w_i\in L$, and matrices
$A_{ji}\in\Mat_3(K)$ of rank at most one, satisfying
\begin{equation}\label{eq:reconstruction}
 \sum_{i=1}^{8}\tr(XA_{ji})w_i=\tr(X)b_j
            \qquad\text{for all }X\in\Mat_3(K),\quad 1\leqs j\leqs5.
\end{equation}
We call the system \eqref{eq:reconstruction} a \emph{reconstruction of
the outputs} $b_1,\ldots,b_5$.
\end{prop}
\begin{proof}
By Lemma~\ref{lem:atom}, the four supports are disjoint and contain
twelve indices. Let $J$ be the eight remaining indices, and let $L$
be the annihilator of the four tests. These tests satisfy the
assumptions of Lemma~\ref{lem:annihilator}; choose its rank-one basis
$b_j$. For $i\in J$, we have $W_i\in L$; set $w_i=W_i$ and
$A_{ji}=V_i b_j U_i$. These matrices have rank at most one.
In \eqref{eq:dual-sandwich} with $T=b_j$, the atomic summands vanish:
if $U_i=ab_i$ and $V_i=c_i s$, then
$V_i b_j U_i=c_i(s b_j a)b_i=0$, since $s b_j a=\tr(b_j as)=0$.
What remains is \eqref{eq:reconstruction}.
\end{proof}

\subsection{Coordinates of a three-dimensional kernel}

\begin{lemma}\label{lem:kernel-model}
If \eqref{eq:reconstruction} holds, then there exist eight pairwise
nonproportional nonzero columns $\mu_i\in K^3$ spanning $K^3$ and,
for each $j$, a linear map $Z_j:K^3\to\Mat_3(K)$ such that
\[
 A_{ji}=a_{ji}I+Z_j(\mu_i),\qquad (Z_j(x))_{11}=0.
\]
Moreover, $I\in\Span_i A_{ji}$, and
\[
 3\leqs\dim\Span_i A_{ji}=1+\rk Z_j\leqs4.
\]
\end{lemma}
\begin{proof}
Set $F:K^8\to L$, $F(c)=\sum_i c_iw_i$.
Substituting $X=E_{11}$ in \eqref{eq:reconstruction} shows that each
$b_j$ lies in the image of $F$. Thus $F$ is surjective and
$\dim\ker F=3$. Let $\mu_i$ be the coordinates of the restriction
of the $i$th coordinate form to $\ker F$, in a chosen basis of
$(\ker F)^*$. They span this entire dual space.

If $\mu_i,\mu_k$ are dependent, a nonzero form
$a\varepsilon_i+b\varepsilon_k$ annihilates $\ker F$ and thus has
form $\ell\circ F$ for some $\ell\in L^*$.
Applying $\ell$ to \eqref{eq:reconstruction} gives
\[
 aA_{ji}+bA_{jk}=\ell(b_j)I.
\]
The left-hand side has rank at most two, so $\ell(b_j)=0$ for every
$j$. Hence $\ell=0$ and $a=b=0$, a contradiction.
This also excludes zero $\mu_i$.

Set $a_{ji}=(A_{ji})_{11}$. For each matrix entry, the vector of
corresponding entries of $A_{ji}-a_{ji}I$ lies in $\ker F$, by
\eqref{eq:reconstruction}. Thus these matrices arise from the
coordinates $\mu_i$ through a linear map $Z_j$; their $(1,1)$ entry
is zero.

Choose $\ell$ with $\ell(b_j)=1$ and apply it to the reconstruction
again: $I=\sum_i\ell(w_i)A_{ji}$. Hence
$\Span_i A_{ji}=KI\oplus\im Z_j$.
This space has dimension at least three: in a basis of matrices of
rank at most one, expressing the rank-three identity requires at
least three matrices. The upper bound of four follows from
$\dim K^3=3$.
\end{proof}

\subsection{Quadratic separation and concurrent lines}

\begin{definition}
A nonzero column $\mu\in K^3$ represents a projective point $[\mu]$;
proportional columns represent the same point. A projective line is
the set of one-dimensional subspaces of a plane in $K^3$.
Three distinct lines are \emph{concurrent} if they pass through a
single point, called their center. For a covering of a family of
points by such lines, the \emph{active group} of a line consists of
the indices of points on it other than the center.

An index $i$ is \emph{quadratically separable} in the family
$[\mu_1],\ldots,[\mu_8]$ if there is a homogeneous quadratic
polynomial $q$ in three variables with $q(\mu_i)\ne0$ and
$q(\mu_k)=0$ for $k\ne i$.
\end{definition}

\begin{lemma}\label{lem:separator}
Let the rank-one matrices $A_i$ be pairwise nonproportional, with
$I\in\Span A_i$, and $\dim\Span A_i\leqs4$.
Then there are an index $i$ and a linear form $\chi$ on matrices such that
\[
 \chi(I)=1,\quad\chi(A_i)\ne0,\quad\chi(A_k)=0\ (k\ne i).
\]
There is also a homogeneous quadratic form $q$ on matrices, invariant
under addition of a scalar multiple of the identity, that separates
the same index.
\end{lemma}
\begin{proof}
Partition the family $A_i$ into connected components; their spans
form a direct sum. Decompose $I=\sum_H T_H$ along this sum and set
$d_H=\dim\Span A_H$, $r_H=\rk T_H$.
Lemma~\ref{lem:connected} gives $2r_H\leqs d_H+1$: the columns of
$T_H$ lie in the span of the first factors of $A_i$, $i\in H$, and
its rows in the span of the second factors, so both spans have
dimension at least $r_H$.
If $d_H\geqs2$, integrality gives $3r_H\leqs2d_H$.
If $T_H=0$ whenever $d_H=1$, we would obtain
\[
 3=\rk I\leqs\sum_H r_H\leqs\frac23\sum_H d_H\leqs\frac83,
\]
a contradiction. Hence a one-dimensional component has $T_H\ne0$.
Pairwise nonproportionality implies that it contains a single $A_i$.
Projection onto this component, extended to a linear form on all
matrices and normalized by its value on $I$, gives $\chi$.

The matrix $A_i-\chi(A_i)I$ has rank at least two: otherwise
$\rk A_i=1$ would imply $\rk I\leqs2$.
Choose a nonzero two-by-two minor of this matrix. Taking that minor
of $T-\chi(T)I$ defines a quadratic form $q(T)$.
For $k\ne i$, it is the corresponding minor of $A_k$, hence zero.
Also $q(T+cI)=q(T)$ because $\chi(I)=1$.
\end{proof}

\begin{lemma}\label{lem:slice-dichotomy}
In the model of Lemma~\ref{lem:kernel-model}, each $j$ satisfies at
least one of the following alternatives:
\begin{enumerate}
\item there is a quadratically separable index $i$ with $b_j\in Kw_i$;
\item the eight points $[\mu_i]$ are covered by three concurrent
lines, each containing at least two points other than the center;
moreover, $b_j$ lies in the span of the $w_i$ in each of the three
active groups.
\end{enumerate}
\end{lemma}
\begin{proof}
Fix $j$ and suppress it in the notation.
If $\dim\Span A_i=4$, the map $Z$ is injective.
Thus the $A_i$ are nonzero and pairwise nonproportional: removing
the scalar part determined by the $(1,1)$ entry from such a relation
would give the same relation among the $\mu_i$.
Apply Lemma~\ref{lem:separator}. The reconstruction and $\chi(I)=1$ give
\[
 b=\sum_k\chi(A_k)w_k=\chi(A_i)w_i.
\]
The quadratic form $x\mapsto q(Z(x))$ separates $\mu_i$, since
\[
 q(Z(\mu_k))=q(A_k-a_kI)=q(A_k).
\]

Now suppose $\dim\Span A_i=3$. Choose a basis of this space consisting
of rank-one matrices. A linear combination of them equals $I$, so
both sets of factors are bases and, after absorbing the nonzero
coefficients, are dual. Hence a conjugation simultaneously
diagonalizes every $A_i$; each nonzero $A_i$ has exactly one nonzero
diagonal coordinate.

After this conjugation, consider
\[
 \pi(x)=\bigl((Z(x))_{22}-(Z(x))_{11},\
                 (Z(x))_{33}-(Z(x))_{11}\bigr).
\]
Its kernel is $\ker Z$: equality of the three diagonal entries means
that the matrix is scalar, while the original condition
$(Z(x))_{11}=0$ excludes a nonzero scalar matrix.
Thus $\rk\pi=2$ and $\dim\ker\pi=1$.
Modulo $KI$, the three coordinate matrix lines give three distinct
lines in $K^2$. Their inverse images under $\pi$ are three planes
through $\ker\pi$, or three concurrent projective lines.
They cover all $[\mu_i]$.

Denote the $a$th diagonal coordinate of the conjugated $A_i$ by $c_{ai}$.
Applying the diagonal coordinate form to the reconstruction gives
$b=\sum_i c_{ai}w_i$ for every $a=1,2,3$.
Nonzero $c_{ai}$ occur on the corresponding line away from the center.
Each active group is nonempty. If all three groups have at least two
indices, we obtain the second alternative. Otherwise one group
consists of a single index $i$, and $b\in Kw_i$.
The product of linear equations for the other two lines is a
quadratic polynomial nonzero at $\mu_i$ and zero at every other point.
This gives the first alternative.
\end{proof}

\begin{lemma}\label{lem:projective-count}
For eight distinct projective points, the following hold.
\begin{enumerate}
\item At most three indices are quadratically separable.
\item If the points are covered by three concurrent lines, each
containing at least two points other than the center, then no index
is quadratically separable.
\item There are at most four distinct centers of such coverings.
\end{enumerate}
\end{lemma}
\begin{proof}
\emph{1.} The evaluation vectors of quadratic forms constitute a
subspace of $K^8$ of dimension at most six. If there are $s$ separable
indices, their coordinate vectors lie in this space. On the remaining
points, evaluation of quadratic forms has rank at least
$\min(3,8-s)$: any three distinct projective points can be separated
by products of linear forms. Thus $s+\min(3,8-s)\leqs6$, which implies
$s\leqs3$.

\emph{2.} Suppose a quadratic form vanishes at all but possibly one
point. Its restriction to a projective line, if nonzero, has at most
two zeros. Thus three zeros on a line force its linear equation to
be a factor of the form.
If the center is not among the eight points, deleting one point
leaves seven points on the three lines, in nonempty groups, at least
two of which have size at least two. We can choose three, two, and
one points on the respective lines. The first triple forces one
linear factor, and the second pair forces the other. The last point,
outside these two lines, gives a contradiction.
If the center is among the eight points and is the point deleted,
seven noncentral points remain, at least two on each line, and the
same argument applies. If a noncentral point is deleted, the center
remains a zero. Two lines still contain at least two noncentral points
each; together with the center these give two triples. The form must
be the product of equations of these two lines, but a remaining
noncentral point on the third line is not a zero of that product.

\emph{3.} If some line contains at least four points, it must occur
in every covering by three lines: otherwise the covering lines
would meet it in at most three points. All centers lie on that line.
For two distinct centers, at most four points can lie off it, at the
four intersections of the remaining pairs of lines; the active-group
condition requires at least four.
Fix a point off the line. In each covering, its covering line contains
another point off the line, with only three possible partners.
A partner determines the covering line and its intersection with the
original line, hence the center. There are therefore at most three centers.

It remains to consider the case where each line contains at most
three points. Every covering has at least two lines containing three
points. Two distinct centers cannot use the same such line: otherwise
the five points off it would lie in the four intersections of the
remaining pairs of lines. Five centers would therefore give ten
distinct three-point lines. Each contains three unordered pairs of
points, and no pair is repeated. But $10\cdot3>\binom82=28$.
\end{proof}

\begin{lemma}\label{lem:same-center}
In the model \eqref{eq:reconstruction}, two independent outputs $b,b'$
cannot have coverings of the second type in
Lemma~\ref{lem:slice-dichotomy} with the same center.
\end{lemma}
\begin{proof}
For a fixed center, each noncentral point determines its line
uniquely. Hence the active groups of the two coverings agree up to
permutation. Let $E_1,E_2,E_3\subseteq K^8$ be the coordinate spaces
of these three disjoint groups. If
$H=F(E_1)\cap F(E_2)\cap F(E_3)$, then $2\dim H\leqs\dim\ker F=3$.
Indeed, choose linear maps $s_a:H\to E_a$ with $F(s_a(h))=h$.
The map $(h,k)\mapsto s_1(h)-s_2(h)+s_1(k)-s_3(k)$ is injective
and takes values in $\ker F$. Injectivity follows by looking at the
second and third disjoint coordinate groups.
Thus $\dim H\leqs1$, and $H$ cannot contain two independent outputs.
\end{proof}

\begin{theorem}\label{thm:no-eight}
Over any field, a reconstruction \eqref{eq:reconstruction} of five
independent outputs by eight terms of rank at most one is impossible.
\end{theorem}
\begin{proof}
Construct the common family of eight points using
Lemma~\ref{lem:kernel-model}. If all five outputs have the first type
in Lemma~\ref{lem:slice-dichotomy}, then by part 1 of
Lemma~\ref{lem:projective-count} they lie in the span of at most three
vectors $w_i$, contradicting independence.
Thus there is at least one covering of the second type.
By part 2, there are now no quadratically separable indices, so all
five outputs have the second type. Their centers are distinct by
Lemma~\ref{lem:same-center}, contradicting part 3 of
Lemma~\ref{lem:projective-count}.
\end{proof}

\begin{cor}\label{cor:no-four}
An algorithm of length twenty cannot have four atomic tests with
distinct supports.
\end{cor}
\begin{proof}
Proposition~\ref{prop:residual-model} reduces such an algorithm to a
reconstruction ruled out by Theorem~\ref{thm:no-eight}.
\end{proof}

\section{Three atoms and completion of the proof}

\begin{prop}\label{prop:three}
For a normalized algorithm of length twenty, the equalities $MD=D$
and $\dim D=3$ are impossible.
\end{prop}
\begin{proof}
By Lemma~\ref{lem:component-cost}, $D$ is spanned by common eigenvectors
for multiplication by $M$. Choose three independent such vectors.
Lemma~\ref{lem:stationary} turns them into atomic tests $C_j=a_js_j$
with distinct supports. By Lemma~\ref{lem:atom}, the columns $a_j$
form a basis, as do the rows $s_j$. Let $(u_j)$ be the row basis dual
to $(a_j)$.

Since $\gamma(C_j)$ is a common eigenvector, $C_{\gamma(C_j)}(u)$ is
a nonzero constant multiple of $(ua_j)C_j$.
Identity \eqref{eq:compatibility} gives
\[
 (ua_j)s_j B_z(v)=(va_j)s_j B_z(u).
\]
Taking $u=u_k$, we get $s_jB_z(u_k)=0$ for $j\ne k$.
Consequently,
\[
 B_z(u_k)=s_k^*f_k(z)
\]
for certain linear maps $f_k:Z\to F$.
Each $f_k$ is injective: a zero would imply $B_z(u_k)=0$, while
$u_k\ne0$. Lemma~\ref{lem:injective-product} gives $r\geqs9z+3$.
But \eqref{eq:dz} with $d=3$ and $r=20$ gives $z\geqs2$, hence
$r\geqs21$, a contradiction.
\end{proof}

\begin{proof}[Proof of the main theorem]
Suppose an algorithm of length $r\leqs20$ exists over the original
field. Extend the field to an algebraic closure and delete zero
summands. Proposition~\ref{prop:critical} shows that the remaining
length is twenty.

Among algorithms of this length, choose one with minimal first rank
sum, and then normalize it using Lemma~\ref{lem:generic}.
Proposition~\ref{prop:no-growth} gives $MD=D$, and
Proposition~\ref{prop:critical} gives $d\geqs3$.
If $d\geqs4$, choose four independent common eigenvectors in $D$.
By Lemmas~\ref{lem:component-cost} and~\ref{lem:stationary}, they give
four atomic tests. Their supports are distinct, since atoms with the
same support are proportional. This contradicts Corollary~\ref{cor:no-four}.
Hence $d=3$, but Proposition~\ref{prop:three} excludes this case.

Thus no algorithm of length at most twenty exists over the original
field either. Therefore $R_K\langle3,3,3\rangle\geqs21$.
\end{proof}
\end{document}
